%% kapa: exam=1 round=21 year=2010 period=1 form=A subject=economics subject_ko=경제원론 seq=28 number=68
%% source: Q-Net 감정평가사 기출문제 (cst00309) artlSeq=5024181 file=extracted/1008/5024181/2010년 제21회 감정평가사 1차시험 기출문제/1교시.hwp
%% conversion: kordoc 4.24.51 -> tools/gen1.py (auto); review=proofread-sonnet; flags=none
\begin{question}{68}
A국과 B국이 어떤 무역정책을 선택하는가에 따라 두 국가의 보수는 다음 표와 같이 결정된다. 보수를 극대화하려는 두 국가가 취할 전략에 관한 설명으로 옳은 것을 모두 고른 것은? (단, 각 보수쌍에서 왼쪽은 A국의 보수이고, 오른쪽은 B국의 보수이다.)

{\small\centering
\begin{tabular}{cc|c|c|}
\cline{3-4}
 & & \multicolumn{2}{c|}{B국의 선택} \\
\cline{3-4}
 & & 자유무역 & 보호무역 \\
\hline
\multicolumn{1}{|c|}{\multirow{2}{*}{A국의 선택}} & 자유무역 & (100, 100) & (-100, 200) \\
\cline{2-4}
\multicolumn{1}{|c|}{} & 보호무역 & (200, -100) & (-50, -50) \\
\hline
\end{tabular}\par}

\begin{framedtext}
ㄱ. 내쉬(Nash) 균형은 두 국가가 모두 보호무역을 선택하는 것이다. \par
ㄴ. 두 국가가 모두 자유무역협정(FTA)을 시행하는 것은 파레토 효율적이다. \par
ㄷ. A국이 보호무역을 선택하는 것은 우월 전략이다. \par
ㄹ. 두 국가가 보호무역 상태에서 어느 한 국가가 무역정책을 변경하면 두 국가 모두 보수가 커진다.
\end{framedtext}
\begin{choices}
  \item ㄱ, ㄴ
  \item ㄱ, ㄷ
  \item ㄷ, ㄹ
  \item ㄱ, ㄴ, ㄷ
  \item ㄴ, ㄷ, ㄹ
\end{choices}
\end{question}
\answer{4}
